\section{Conclusion}
\label{sec:conclusion}

Cartel enforcement often begins before the agency has the
evidence it would need to prove a case. Procurement records then
pose a problem that is prior to liability: how should scarce
investigative capacity be allocated when the cheap administrative
record is broad but legally thin, and the richer bid-level record
is costly to recover and evaluate? This paper answers that
question by treating cartel screening as evidence allocation
under costly proof.

The empirical design separates the award layer from the bid
layer. The award layer observes participation and outcomes; it
can rank persistent zero-win participation before bid histories
are opened. The bid layer observes submitted offers and
within-tender patterns; it is the natural place for
proof-producing forensic evaluation. The contribution is not to
collapse these layers. It is to show how they can be sequenced:
administrative records identify where investigation should begin,
and bid-level evidence then evaluates conduct inside the
prioritized set.

The results support that architecture. Persistent zero-win
participation is enriched among adjudication-anchored
cartel-adjacent losers, survives timing and leakage discipline,
and does not operate as a generic direct-defendant detector. The
cobidder profile gives the target economic content: within the
frequent-loser stratum, cartel-adjacent firms are more broadly
deployed, closer to legal anchors, more product-concentrated, and
distinct in bid-level behavior. The bid-distribution exercises
then show complementarity rather than substitution. Award-layer
ranking adds information to bid-layer moments, and a sequential
gatekeeper preserves much of the adjudication signal while
reducing the bid-microdata footprint by
$\valArchFootprintReduction$.

The price evidence reinforces the same legal-economic posture.
Flagged environments have a positive broad price imprint, but
overlap and segment results prevent that imprint from becoming a
damages or overcharge claim. The screen is valuable because it
orders forensic priority, not because it decides harm. Its proper
use is institutional and staged: rank where costly evidence
should be assembled, then let bid records, documents, and legal
process determine what the evidence proves.

That staged use is the broader lesson. Agencies routinely face a
gap between what they can observe cheaply and what they must know
to prove coordination. Waiting for full proof before prioritizing
investigation wastes the cheap information already in the record;
treating cheap suspicion as proof overstates what the record can
bear. The approach in this paper occupies the middle ground. It
turns a thin administrative signal into a disciplined queue for
costly legal-evidentiary work.
